\section{Regla de Oro de Seguridad}

El Panel Gerencial tiene un inmenso nivel de poder. Un clic equivocado puede desbalancear toda la historia médica de la clínica, por lo cual es vital familiarizarse con las mejores prácticas preventivas.

\subsection{Prohibición Absoluta: Nunca uses "Delete"}

Al navegar por usuarios, pacientes, o estudios, siempre verás botones y selecciones con la palabra gigantesca \textbf{Delete} (Eliminar).

Como administrador de salud, \textbf{jamás debes presionar "Eliminar" un registro consolidado (usuarios o estudios del paciente)}. Esto destruiría evidencia de facturación, borraría dictaminaciones pasadas del historial global y rompería las leyes de integridad de datos médicos (conservación de la historia clínica).

Sólo puedes eliminar registros si estás 100\% seguro de que fue un error tipográfico inservible (ej. alguien entró a probar que los botones servían, o llenó campos con el nombre "Prueba Prueba").

\subsection{La Solución Legal: Suspender o "Desactivar"}

Si sorprendes que la recepcionista actual está actuando incorrectamente, roba fondos, decide renunciar sorpresivamente, o es despedida, NO LA BORRES ni borres a ningún médico en su primer mes de renuncia. Haz lo siguiente para retirarle instantáneamente el poder, conservando su registro:

\begin{enumerate}
    \item Ve rápidamente a \textbf{Users} en el panel general.
    \item Encuentra el correo de ese doctor, paciente o asistente con problema. Trata de abrirlo.
    \item Busca la pequeña casilla que dice \textbf{"Active"} (Activo).
    \item \textbf{Desmarca esa casilla (quita la palomita).}
    \item Presiona Save (Guardar) hasta abajo.
    \item A partir de ese segundo crítico, su credencial queda virtualmente congelada y será rechazada de intentar loguearse a tu clínica desde casa o computadora ajena, resolviendo el problema sin quebrar las referencias y enlaces de los documentos que ya emitió como trabajador.
\end{enumerate}

\begin{figure}[H]
    \centering
    \includegraphics[width=0.55\linewidth]{screenshots/suspend_account.png}
    \caption{El Gerente desactivando sabiamente a un ex-empleado en vez de borrarlo.}
\end{figure}

\subsection{Auditoría de Supervisores ("Staff Status")}

Te recomendamos calendarizarte al último día de cada mes entrar a la bandeja de "Usuarios" (Users) en este panel directivo para revisar quienes poseen la palomita habilitada llamada \textbf{Staff Status} (Estatus Institucional del Staff). Quien posea esa palomita, tiene acceso libre a este propio portal que tú estás visualizando en estos momentos. Asegúrate de que no existan ex-empleados infiltrados con esa viñeta superior puesta.
